% Part: methods
% Chapter: proofs
% Section: proof-by-contradiction

\documentclass[../../../include/open-logic-section]{subfiles}

\begin{document}

\olfileid{mth}{prf}{con}

\olsection{د تناقض له لارې ثبوت}

په لومړي سر کښې د تناقض له لارې ثبوت د استنتاج يوه بڼه
ده چې منفي دعوې پرې ثابتېږي۔ فرض کړئ غواړئ وښيئ چې کومه
دعوه~$p$ \emph{ناسمه} ده؛ يعنې غواړئ~$\lnot p$ ثابت کړئ۔
تر ټولو غوره لاره دا ده: (الف) فرض کړئ چې $p$ رښتيا ده،
او (ب) وښيئ چې له دې فرضه داسې پايله راځي چې ناسموالی
يې درته څرګند دے۔ «د څرګند ناسموالي پايله» ښايي له
خپلې~$p$ سره يا د کتل کېدونکې ټولې دعوې له کوم بل
فرض سره په ټکر کښې وي، يعنې تناقض جوړ کړي۔ د بېلګې
په توګه، د «که $q$ وي نو $\lnot p$» ثبوت کښې $q$
رښتيا فرضېږي او له هغې~$\lnot
p$ ثابتېږي۔ که $\lnot p$ د تناقض له لارې ثابتوئ،
دا د $p$ فرضول دي، د~$q$ ترڅنګ۔ که $\lnot q$
له $p$ څخه ثابت کړے شئ، ښيئ چې د~$p$ له فرضه
ستاسې له بل فرض~$q$ سره په ټکر کښې پايله راځي، ځکه
$q$ او $\lnot q$ دواړه په يو وخت رښتيا کېدلے نه شي۔
البته، د تناقض د ثبوت په بهير کښې نورې استنتاجي بڼې او
د تعريفونو پرانيستل هم کارول پکار دي۔ يوه بېلګه ګورو۔

\begin{prop}
  که $A \subseteq B$ او $B = \emptyset$ وي، نو $A$ هېڅ
  !!{element}s نه لري۔
\end{prop}

\begin{proof}
  فرض کړئ $A \subseteq B$ او $B = \emptyset$۔ غواړو وښيو چې
  $A$ هېڅ !!{element}s نه لري۔
  \begin{quote}
    دا شرطي دعوه ده؛ نو مقدم فرضوو او نتيجه ثابتوو۔ نتيجه
    دا ده چې $A$ هېڅ !!{element}s نه لري۔ په څرګنده بڼه:
    داسې~$x \in A$ نشته۔
  \end{quote}
  $A$ هېڅ !!{element}s نه لري هغه وخت او يوازې هغه وخت
  چې داسې~$x$ نشته چې $x \in A$ وي۔
  \begin{quote}
    نو اصلاً يوه منفي دعوه~$\lnot p$ ثابتول غواړو: داسې
    $x \in A$ نشته۔ د تناقض له لارې ثبوت دپاره يې مقابل
    مثبت حکم~$p$ فرضوو، يعنې داسې $x \in A$ شته، او
    له دې څخه تناقض ثابتوو۔ دا لاره په «د تناقض دپاره
    فرض کړئ» يا لنډ «برعکس فرض کړئ» پېلوو او بيا
    فرض~$p$ ليکو۔
  \end{quote}
  برعکس فرض کړئ: داسې $x \in A$ شته۔
  \begin{quote}
    دا اوس هغه نوے فرض دے چې تناقض ترې اخلو۔ دوه نور
    فرضونه هم لرو: $A \subseteq B$ او $B = \emptyset$۔
    له لومړي څخه $x \in B$ راځي:
  \end{quote}
  څرنګه چې $A \subseteq B$، نو $x \in B$۔
  \begin{quote}
    خو څرنګه چې $B = \emptyset$، د $B$ هر !!{element}،
    لکه $x$، بايد د~$\emptyset$ !!a{element} هم وي۔
  \end{quote}
  څرنګه چې $B = \emptyset$، نو $x \in \emptyset$۔ دا تناقض
  دے، ځکه د تعريف له مخې $\emptyset$ هېڅ !!{element}s
  نه لري۔
  \begin{quote}
    ثبوت همدلته بشپړېږي: له ټاکلو فرضونو څخه تناقض ته
    ورسېدو، نو ټول فرضونه يو ځای رښتيا نه شي کېدلے۔
    لومړي دوه فرضونه، $A \subseteq B$ او $B = \emptyset$،
    د قضیې ورکړل شوي شرطونه دي؛ نو وروستی زيات شوے فرض،
    دا چې داسې $x \in A$ شته، بايد ناسم وي۔ که بشپړ
    تفصيل غواړو، دا پايله ښکاره ليکلے شو۔
  \end{quote}
  نو زموږ دا فرض چې داسې $x \in A$ شته ناسم دے؛ له دې امله
  $A$ د تناقض له لارې د ثبوت له مخې هېڅ !!{element}s نه لري۔
\end{proof}

هره مثبته دعوه په څرګنده توګه له يوې منفي
دعوې سره برابره ده: $p$ هغه وخت او يوازې هغه وخت چې
$\lnot\lnot p$۔ نو د تناقض له لارې ثبوت په نامستقيم ډول د
مثبتو دعوو د ثابتولو دپاره هم کارېدلے شي۔ د~$p$ د ثبوت
دپاره يې د منفي دعوې $\lnot\lnot p$ په توګه واخلئ۔
که له $\lnot p$ څخه تناقض ثابت کړو، د تناقض له لارې
مو $\lnot\lnot p$ ثابت کړ، او له هغې~$p$۔

په تېرې بېلګې کښې مو يوه منفي دعوه ثابتول غوښتل:
دا چې $A$ هېڅ !!{element}s نه لري۔ نو د تناقض د ثبوت
دپاره مو مثبت فرض ومانۀ، يعنې داسې $x \in A$ شته۔
دې موږ ته د کار دپاره فرضي $x \in A$ راکړ؛ ورپسې مو
تر $x \in \emptyset$ پورې د هغه په اړه استدلال وکړ۔

کله چې يوه مثبته دعوه په نامستقيم ډول ثابتوئ، د
تناقض دپاره منفي فرض نيسئ۔ خو ډېر ځله مثبته دعوه په
منفي بڼه او منفي دعوه په مثبته بڼه په اسانۍ اړولے شئ۔
زموږ تېر ثبوت به تقريباً هماغسې وو که مو
«$A = \emptyset$» ثابتول غوښتل، د «$A$ هېڅ
!!{element}s نه لري» د منفي نتيجې پر ځاے۔
(د $=$ له تعريفه «$A = \emptyset$» يوه
عمومي دعوه ده: «د~$A$ هر !!{element} د~$\emptyset$
!!{element} دے او برعکس».) خو په اسانۍ ښکاري چې
دا له «نه: داسې $x \in A$ شته» سره برابره ده۔

نو کله ناکله له $\lnot p$ سره د فرض په توګه کار
کول د~$p$ له نېغ ثبوته اسان وي۔ که نېغ
ثبوت همدومره اسان يا تر دې اسان هم وي، لکه راتلونکو
بېلګو کښې، ځينې خلک نامستقيمه لاره خوښوي۔ که دوه
ځله نفي مو مغشوشوي، د يوې دعوې د تناقض له لارې ثبوت
د هغې د \emph{مقابلې} دعوې له فرضه د تناقض اخيستل وګڼئ۔
نو د $\lnot
p$ د تناقض له لارې ثبوت د~$p$ له فرضه تناقض
ثابتوي؛ او د~$p$ د تناقض له لارې ثبوت له~$\lnot p$
څخه تناقض ثابتوي۔

\begin{prop}
$A \subseteq A \cup B$۔
\end{prop}

\begin{proof}
  غواړو وښيو چې $A \subseteq A \cup B$۔
  \begin{quote}
    په ظاهر کښې دا مثبته دعوه ده: هر $x \in A$ د
    $A \cup B$ غړے هم دے۔ نفي يې دا ده: ځينې
    $x \in A$، $\notin A \cup B$ دي۔ نو دا منفي دعوه فرضولے شو
    او ښودلے شو چې تناقض ته رسېږي؛ په دې توګه دعوه
    په نامستقيم ډول ثابتېږي۔
    \end{quote}
  برعکس فرض کړئ، يعنې $A \nsubseteq A \cup B$۔
  \begin{quote}
    د $A \subseteq A \cup B$ تعريف لرو: هر $x \in A$،
    $\in A \cup B$ هم دے۔ د
    $A \nsubseteq A \cup B$ د معنا د پوهېدو دپاره
    د پرانيستي تعريف ساده منطقي اړونه کاروو: دا خبره
    ناسمه ده چې هر $x \in A$، $\in
    A \cup B$ هم دے، هغه وخت او يوازې هغه
    وخت چې \emph{يو}~$x \in A$، $\notin C$ وي۔
    \emph{د سرچينې يادونه (OLSTH-037): دلته د «سي» نښه
    د روانې دعوې د اتحاد پر ځاے چاپ شوې ده؛ لاندې
    بشپړ شرط هماغه اتحاد کاروي۔ اصلي فورمول ساتل شوے دے۔}
    (په دې ځاے کښې ډېر پام پکار دے:
    د زده کوونکو د تناقض د ثبوت ډېرې هڅې د «ټول
    $A$s د $B$s دي» په څېر دعوې د نفي له ناسمه
    شننه لوېږي۔) په بل عبارت،
    $A \nsubseteq A \cup B$ هغه وخت او يوازې هغه وخت
    چې داسې $x$ وي چې $x \in A$ او
    $x \notin A \cup B$۔ له دې وروسته کار اسان دے۔
  \end{quote}
  نو داسې $x \in A$ شته چې $x \notin A \cup B$۔
  د $\cup$ له تعريفه، $x \in A \cup B$ هغه وخت او
  يوازې هغه وخت چې $x \in A$ يا $x \in
  B$ وي۔ څرنګه چې $x \in A$، نو $x \in A \cup B$۔
  دا له هغه فرض سره په ټکر کښې دے چې
  $x \notin A \cup B$۔
\end{proof}

\begin{prob}
په \emph{نامستقيم ډول} ثابت کړئ چې $A \cap B \subseteq A$۔
\end{prob}

\begin{prop}
که $A \subseteq B$ او $B \subseteq C$ وي، نو $A \subseteq C$۔
\end{prop}

\begin{proof}
  فرض کړئ $A \subseteq B$ او $B \subseteq C$۔ غواړو وښيو چې
  $A
  \subseteq C$۔
  \begin{quote}
    نامستقيمه لاره نيسو: د غوښتل شوې پايلې نفي فرضوو۔
  \end{quote}
  برعکس فرض کړئ، يعنې $A \nsubseteq C$۔
  \begin{quote}
    د تېر په شان، $A \nsubseteq C$ معنا دا ده چې دا خبره
    رښتيا نۀ ده چې هر $x \in A$، $\in C$ هم وي؛
    يعنې ځينې $x \in A$، $\notin C$ دي۔ اندېښنه مه کوئ:
    په تمرين سره به د داسې نفي د پرانيستلو دپاره
    ډېر فکر ته اړ نه شئ۔
  \end{quote}
  په بل عبارت، داسې~$x$ شته چې $x \in A$ او
  $x \notin C$۔
  \begin{quote}
    اوس ترې تناقض اخيستلے شو۔ البته، د دې دپاره
    نور دوه فرضونه هم کاروو۔
  \end{quote}
  څرنګه چې $A \subseteq B$، نو $x \in B$۔ څرنګه چې
  $B \subseteq C$، نو $x \in
  C$۔ خو دا له $x \notin C$ سره تناقض دے۔
\end{proof}

\begin{prop}
که $A \cup B = A \cap B$ وي، نو $A = B$۔
\end{prop}

\begin{proof}
  فرض کړئ $A \cup B = A \cap B$۔ غواړو وښيو چې $A = B$۔
  \begin{quote}
    پيل اوس عادي دے:
  \end{quote}
  د تناقض دپاره فرض کړئ چې $A \neq B$۔
  \begin{quote}
    د تناقض د ثبوت فرض دا دے چې $A \neq
    B$۔ څرنګه چې $A = B$ هغه وخت او يوازې هغه وخت
    چې $A \subseteq B$ او $B \subseteq A$، نو
    $A \neq B$ هغه وخت او يوازې هغه وخت چې
    $A \nsubseteq B$ \emph{يا} $B \nsubseteq
    A$۔ (وګورئ چې د نفيو په اړولو کښې پام څومره اړين
    دے!{}) له دې فصله د تناقض د اخيستلو دپاره د حالتونو
    له مخې ثبوت کاروو او ښيو چې په هر حالت کښې تناقض
    راځي۔
  \end{quote}
  $A \neq B$ هغه وخت او يوازې هغه وخت چې
  $A \nsubseteq B$ يا $B \nsubseteq A$۔ حالتونه بېلوو۔
  \begin{quote}
    په لومړي حالت کښې $A \nsubseteq B$ فرضوو؛ يعنې
    د ځينو $x$ دپاره $x \in A$ خو $\notin B$۔ $A \cap B$ د هغو
    !!{element}s سټ دے چې $A$ او $B$ دواړه يې لري؛
    نو که يو څيز په يوه کښې نه وي، په تقاطع کښې هم
    نه وي۔ $A \cup B$ د $A$ او $B$ يوځای کول دي؛
    نو د دواړو له هر يوه غړے په اتحاد کښې هم وي۔
    له دې اخلو چې $x \in A \cup B$ خو
    $x \notin A \cap
    B$، او په پايله کښې $A \cap B \neq A \cup B$۔
  \end{quote}

  حالت ۱: $A \nsubseteq B$۔ نو د ځينو $x$ دپاره
  $x \in A$ خو $x \notin
  B$۔ څرنګه چې $x \notin B$، نو
  $x \notin A \cap B$۔ څرنګه چې $x \in A$،
  $x \in A \cup B$۔ نو $A \cap B \neq A \cup B$؛
  دا له هغه فرض سره په ټکر کښې دے چې
  $A \cap B = A \cup B$۔

  حالت ۲: $B \nsubseteq A$۔ نو د ځينو $y$ دپاره
  $y \in B$ خو $y \notin
  A$۔ لکه مخکې، $y \in A \cup B$ لرو خو
  $y \notin A \cap B$؛ نو $A \cap B \neq A \cup B$،
  او دا بيا له $A \cap B = A \cup
  B$ سره په ټکر کښې دے۔
\end{proof}

\end{document}
